\documentclass[11pt,a4paper]{article}

\usepackage{fontspec}
\usepackage{polyglossia}
\setdefaultlanguage{russian}
\setotherlanguage{english}
\setmainfont{PT Serif}
\setsansfont{DejaVu Sans}
\setmonofont{DejaVu Sans Mono}[Scale=0.82]

\usepackage[margin=2.5cm]{geometry}
\usepackage{amsmath,amssymb}
\usepackage{booktabs}
\usepackage{longtable}
\usepackage{tabularx}
\usepackage{graphicx}
\usepackage{xcolor}
\usepackage{enumitem}
\usepackage{caption}
\captionsetup{font=small,labelfont=bf}
\usepackage{microtype}
\setlength{\emergencystretch}{3em}
\usepackage{pgfplots}
\pgfplotsset{compat=1.18}
\usetikzlibrary{positioning,arrows.meta}
\usepackage[numbers,sort&compress]{natbib}
\usepackage[hidelinks]{hyperref}
\usepackage{url}

\definecolor{skeinc}{HTML}{2F6F9F}
\definecolor{openc}{HTML}{C66A2B}
\definecolor{freshc}{HTML}{2F6F9F}
\definecolor{cachec}{HTML}{9CB8CC}
\definecolor{outc}{HTML}{C66A2B}

\title{\textbf{Skein: агент разработки с немонотонным контекстом}\\[0.4em]
\large Архитектура и измерения на портированных задачах Terminal-Bench}
\author{Владимир Лапшин\\[0.3em]
\small Независимый исследователь\\[0.2em]
\small \texttt{mefrill@yandex.ru}}
\date{Препринт, ноябрь 2026}

\begin{document}

\maketitle

\begin{abstract}
Агент на языковой модели рассуждает вероятностно: его выводы правдоподобны, но не доказуемы.
Поэтому надёжность достигается не внутри модели, а вокруг неё --- процедурой-дисциплиной. Она
задаёт, что модель обязана записать, на что может опираться и когда работа считается
закрытой. Мы описываем Skein, агента разработки, в котором эта процедура встроена в журнал
работы.

Журнал хранит все действия и наблюдения и не переписывается. Сообщения для модели каждый ход
собираются из журнала заново, поэтому лишнее можно убрать из вида, не потеряв. Поверх журнала
работают два явных способа. Состояния задают цель, условие завершения и правила перехода
заранее. Динамические гипотезы позволяют модели самой выдвигать проверяемые утверждения,
работать под ними и закрывать их, записав, что именно она увидела. В обоих случаях движок
фиксирует, что сделано и на каком основании, а не полагается на пересказ модели.

Мы реализовали агента и проверили его на портированных задачах Terminal-Bench. Состояния и
гипотезы решают задачи одинаково; гипотезы тратят почти вдвое меньше новых входных токенов,
но делают больше ходов. Главное ограничение то же, что у любого такого агента: решение модели
не заменяет проверку. Результат принимает внешняя проверка, и именно она показывает, где
модель ошиблась.
\end{abstract}

\noindent\textbf{Ключевые слова:} агент разработки, агенты на больших языковых моделях,
управление контекстом, журнал событий, событийное проектирование (event sourcing), проекция контекста, изоляция
контекста, кэширование префикса, состояния работы, Terminal-Bench.

\section{Введение}

Агент разработки --- программа, которая читает проект, вносит изменения и проверяет
результат, опираясь на решения языковой модели. Типичная реализация следует циклу
ReAct~\cite{yao2023react}: модель рассуждает, выбирает действие, исполнитель запускает
инструмент и возвращает наблюдение. Историю такого цикла естественно хранить как ленту
сообщений: системная инструкция, запрос пользователя, пары «действие--наблюдение». Лента
монотонна --- она только растёт.

У растущей ленты три известные проблемы. Во-первых, контекст упирается в окно модели на
длинных задачах. Во-вторых, модели неравномерно используют длинный контекст: информация
в середине доступна хуже, чем в начале и конце~\cite{liu2023lost}; отдельно от этого
позиционного эффекта мы наблюдаем инженерный феномен --- исходная цель «тонет» в глубине
истории (далее --- дрейф цели; это наблюдение, а не результат~\cite{liu2023lost}).
В-третьих, сжатие и усечение либо теряют факты, либо переписывают их, и точное состояние
работы становится невосстановимым.

Skein строится на другом допущении: \emph{история --- это данные, а контекст ---
производное представление}. Движок записывает действия и наблюдения в дописываемый
типизированный журнал событий (в коде --- IR; далее --- журнал); сообщения для модели на каждом ходу собираются из журнала функцией
$project()$. Такое разделение --- прямой аналог событийного проектирования (event sourcing)
и разделения моделей записи и чтения (CQRS)~\cite{fowler2005event,fowler2011cqrs}:
журнал --- единственная записываемая модель, контекст --- пересчитываемая модель чтения.
Проекция детерминирована как чистая функция и немонотонна как вид\footnote{Здесь
«немонотонный» означает, что видимое представление может сокращаться при смене фокуса; это
не немонотонная логика, где выводы могут отзываться.}: содержимое уходит из
контекста при смене фокуса, не исчезая из журнала.

Второе решение --- явно задавать, из каких шагов состоит работа, и хранить это в журнале, а не
в коде. Здесь есть два способа. В состояниях шаги (фазы) заданы заранее: у каждого есть цель и
условие завершения, а правила перехода описаны протоколом; решение о переходе принимает
модель, исполнитель записывает переход в журнал. В динамических гипотезах шаги задаёт сама
модель по ходу работы: она выдвигает проверяемое утверждение, работает под ним и закрывает
его, записав, что увидела. В обоих случаях цель, условие проверки и факты работы хранятся в
журнале, поэтому к ним можно вернуться и проверить их после работы. Этим Skein отличается и
от одноразовых субагентов, и от графов, зашитых в код (раздел~\ref{sec:states}).

Вклад работы:
\begin{enumerate}[leftmargin=1.4em,itemsep=0.2em]
  \item архитектура, в которой контекст агента --- детерминированная немонотонная проекция
  дописываемого журнала, с адресуемым доступом к скрытой истории (\texttt{recall},
  \texttt{search}); позиционирование относительно event sourcing, OpenHands, MemGPT и
  субагентов;
  \item протокол состояний: цель, критерий остановки и топология переходов как данные
  журнала, интерпретируемые протоколом, а не кодом;
  \item контролируемое сравнение с opencode на 13 портированных задачах Terminal-Bench
  ($K=5$, два сетевых режима) с декомпозицией токенов и стоимости по состояниям;
  анализ структуры затрат и роли кэширования префикса в экономике длинной ленты;
  \item динамические гипотезы --- следующий шаг от явно заданных состояний: вместо фаз модель
  формулирует и проверяет утверждения; сравнение с состояниями на семи задачах ($K=5$, два
  сетевых режима), включая длинную задачу и задачу, где программу пишут с нуля;
  \item разбор провалов и план, как измерить то, что пока не измерено: удерживает ли агент
  цель и насколько им можно управлять.
\end{enumerate}

Код, документация и инструкция по установке доступны в репозитории
\url{https://github.com/merfill/skein}, на сайте \url{https://skein-swe.dev} и как пакет
npm \texttt{skein-swe} (\url{https://www.npmjs.com/package/skein-swe}).

Работа устроена вокруг одного вопроса: как организовать рассуждение модели так, чтобы его
можно было проверить после того, как оно сделано. Ядро Skein небольшое: дописываемый журнал,
проекция в сообщения и два способа работы поверх журнала --- состояния и динамические
гипотезы. Этого достаточно, чтобы измерить цену разделения контекста и показать, как устроена
работа с явными шагами, а не со свободной лентой.

С гипотезами связано и более общее предположение. Сейчас модели учат в основном прямым
рассуждениям --- цепочке «подумал и сразу ответил», --- и такая цепочка часто непредсказуема.
Гипотезы задают другую рамку: сначала выдвинуть утверждение, потом проверить его и только
потом двигаться дальше. В этом смысле они метаорганизатор рассуждения. Мы не утверждаем, что
Skein уже даёт такую организацию в готовом виде; мы показываем, как она устроена, и надеемся,
что такие подходы подтолкнут к обучению моделей рассуждению через гипотезы, а не только через
прямые цепочки.

Не-цели работы: заявление о лучшем результате на наборе задач, оптимизация стоимости,
проверка переносимости между моделями. При одном поставщике API и одной модели сравнение
показывает, что Skein не дешевле opencode; цель статьи --- описать механику и измерить её
последствия.

\section{Связанные работы}

\subsection{Агенты разработки и фазовые пайплайны}

SWE-bench~\cite{jimenez2024swebench} сделал стандартом проверку агентов на реальных задачах
из репозиториев; Terminal-Bench~\cite{terminalbench} переносит идею в терминал. SWE-agent
показывает, что дизайн интерфейса «агент--компьютер» существенно влияет на
результативность~\cite{yang2024sweagent}. Agentless демонстрирует, что жёсткий фазовый
пайплайн «локализация--исправление--валидация» без автономного планирования конкурентен с
агентами~\cite{xia2024agentless}. Последовательность состояний Skein близка к такому
пайплайну по форме, но переходы между фазами выбирает модель по протоколу, а не
расписание в коде; цель и критерий каждой фазы --- данные журнала.

opencode~\cite{opencode} --- открытый терминальный агент, наш эталон сравнения: одна
растущая лента, набор инструментов (оболочка, чтение, поиск, правка, запись), цикл до
остановки модели или исчерпания лимита. Это представитель основного класса реализаций.

\subsection{История как данные: журнал, проекция, память}

Событийное проектирование (event sourcing)~\cite{fowler2005event} и разделение моделей записи и чтения (CQRS)~\cite{fowler2011cqrs} задают схему:
состояние приложения полностью выводится из журнала событий, модели записи и чтения
разделены. Skein применяет эту схему к контексту агента на большой языковой модели: журнал ---
источник истины, контекст --- производное представление. Наш вклад --- не схема, а её
применение и контракт состояний поверх него.

Lost in the Middle~\cite{liu2023lost} мотивирует отказ от длинной ленты, но измеряет
позиционную доступность, а не дрейф цели. MemGPT~\cite{packer2023memgpt} управляет памятью
иерархически, и его основной контекст тоже немонотонен. Ключевое отличие Skein --- в
способе изменения вида: MemGPT \emph{редактирует} представление модельными операциями
памяти, Skein \emph{пересчитывает} представление из неизменяемого журнала. Редактирование
накапливает решения модели о том, что помнить; пересчёт --- нет. CoALA~\cite{sumers2024coala} даёт таксономию памяти и пространств действий; Generative
Agents~\cite{park2023generative} использует поток памяти с извлечением --- редактируемую
память, а не пересчёт из неизменяемого журнала; Letta~\cite{letta} --- продуктовая версия
линии MemGPT. OpenHands~\cite{openhands2024,openhandsSdk2025} --- ближайший родственник:
событийный поток и Condenser, который формирует представление модели, отбрасывая события и
заменяя их сводками. Конденсация OpenHands редактирует и сжимает представление (линия
MemGPT); Skein детерминированно пересчитывает его из неизменяемого журнала и сохраняет
адресуемый доступ к скрытому. Наш вклад --- не event sourcing сам по себе, а его применение
к контексту агента вместе с контрактом состояний в журнале. Reflexion добавляет вербальную обратную
связь поверх ленты~\cite{shinn2023reflexion}. LangGraph
персистентно сохраняет состояние графа чекпойнтами~\cite{langgraphPersistence}, но и
топология, и контракты узлов там задаются кодом.

\subsection{Субагенты и изоляция контекста}

Схема «координатор--исполнители» описана в обзоре Anthropic~\cite{anthropic2024building} и
в их многоагентной исследовательской системе~\cite{anthropic2025multiagent}. В deepagents
субагент --- явная сущность: инструмент \texttt{task}, изолированный контекст, собственная
системная инструкция и набор инструментов~\cite{langchain2025deepagents}; мотивация ---
не засорять контекст родителя промежуточными результатами~\cite{dbreunig2025context}.

Функционально операции Skein \texttt{general}, \texttt{revision}, \texttt{plan},
\texttt{write\_document} --- субагенты в этом смысле. Отличия механики: (i) изоляция
структурная --- у каждой единицы работы свой журнал и своя проекция, а не отдельная
лента; (ii) у состояний внутри \texttt{task} цель и критерий остановки --- данные журнала,
а топология переходов описывается протоколом. В deepagents есть \texttt{CompiledSubAgent}
--- произвольный LangGraph-граф, технически тоже автомат; но его цель и критерий заданы
кодом графа, а не данными, и потому не переносимы между топологиями без изменения кода.
Спецификация субагента в deepagents тоже декларативна, поэтому новизна не в декларативности
как таковой, а в том, что контракт и факты переходов живут в журнале и проверяемы
постфактум. Декларативные фреймворки вроде CrewAI~\cite{crewai} задают роли и задачи в YAML,
но не журналируют переходы и не пересчитывают контекст из неизменяемого журнала. Это
Существенно здесь не «данные против кода» само по себе, а то, что контракт и факты
переходов живут в журнале и проверяемы постфактум; это мы и обсуждаем в разделе~\ref{sec:states}.

\section{Архитектура}
\label{sec:arch}

\subsection{Термины}

Ниже --- термины, которыми мы пользуемся дальше.

\begin{table}[h]
\centering
\small
\caption{Термины и их отношения.}
\label{tab:terms}
\begin{tabularx}{\textwidth}{@{}lX@{}}
\toprule
Термин & Определение \\
\midrule
Сессия & упорядоченный набор единиц работы; агент сессии маршрутизирует сообщения \\
Единица работы & изолированный журнал + проекция; открывается операцией \\
Операция & тип единицы работы (\texttt{respond}, \texttt{general}, \ldots; табл.~\ref{tab:ops}) \\
Состояние & узел протокола внутри единицы работы: цель, критерий, инструменты, переходы \\
Журнал & дописываемая последовательность типизированных событий и узлов (в коде --- IR) \\
Проекция & чистая функция $project(journal) \rightarrow \texttt{Context}$ \\
Ход & один вызов модели (цикл «сообщение $\rightarrow$ действие $\rightarrow$ наблюдение»);
у opencode учтён шаг агента \\
\bottomrule
\end{tabularx}
\end{table}

\subsection{Журнал и проекция}

Журнал монотонен: события только добавляются, узлы не переписываются, полная история
остаётся адресуемой, даже когда не показана модели. Узел --- запись с полями \texttt{id},
\texttt{space}, \texttt{kind}, \texttt{label}, \texttt{payload}, \texttt{seq}; ребро ---
\texttt{id}, \texttt{from}, \texttt{to}, \texttt{kind}, \texttt{provenance}. События:
\texttt{add\_node}, \texttt{add\_edge}, \texttt{descend}, \texttt{return},
\texttt{mutate}, \texttt{revert}, \texttt{rollback}, \texttt{record\_rejection}. Типы
узлов: \texttt{request}, \texttt{move}, \texttt{turns}, \texttt{action},
\texttt{observation}, \texttt{stop}, \texttt{unactionable}, \texttt{constraint},
\texttt{fact}, \texttt{hint}, \texttt{memo}, \texttt{plan}, \texttt{plan\_item},
\texttt{file}, \texttt{symbol}, \texttt{test}. Типы рёбер: \texttt{move},
\texttt{unactionable}, \texttt{turns}, \texttt{turn}, \texttt{stop}, \texttt{result},
\texttt{mutates}, \texttt{plan}, \texttt{plan\_item}. События \texttt{mutate},
\texttt{revert} и \texttt{rollback} не переписывают узлы, а фиксируют новую версию
артефакта или откат; версии хранятся в журнале. «Заметка модели» в рабочей инструкции ---
это поле \texttt{what} события \texttt{move}: оно записывается при входе в состояние и не
переписывается; проекция читает его как непроверенную догадку.

Проекция детерминирована: контекст полностью определяется журналом и параметрами
проекции, что делает возможным воспроизведение контекста по журналу. Подчеркнём:
детерминированность проекции \emph{не} означает детерминированности поведения агента ---
решения принимает стохастическая модель. Свойство проверено на 400 случайных журналах
(\texttt{tests/ops/ir\_properties.test.ts}) и на 125 журналах кампании, сохранённых в
\texttt{result.json} клеток: два вызова \texttt{project()} на одном журнале дают одинаковый
по содержимому контекст (125/125).

Проекция немонотонна. Открытое состояние показывает:
\begin{itemize}[leftmargin=1.4em,itemsep=0.15em]
  \item \texttt{Start state} и краткое описание;
  \item свои ходы --- монотонный список действий и наблюдений;
  \item рабочую инструкцию хода: \texttt{Goal} (цель состояния), непроверенную заметку
  модели, \texttt{Done when} (критерий остановки), \texttt{Plan} (инструкция состояния) и
  явную пометку, что исходный запрос --- контекст, а не цель состояния.
\end{itemize}
Закрытое состояние показывает только \texttt{Start state} и \texttt{Finish state}.
Закрытие не удаляет узлы: они доступны через \texttt{recall} и \texttt{search}.
\texttt{recall(id, start?, end?)} читает сохранённый результат по идентификатору из ленты
(без окна --- с начала, с окном --- страницы большого тела); \texttt{search(id, pattern,
before?, after?)} ищет регулярное выражение внутри сохранённого результата (stdout и stderr)
и возвращает окна строк с контекстом. Оба инструмента только читают, доступны всегда и
добавляют возвращённый текст в ленту открытого состояния; повторный \texttt{recall} без
окна, когда тело уже в виде, отклоняется. Отдельного бюджета токенов нет: размер задаётся
окном и пагинацией. Вид может уменьшаться, не теряя данных.

Сообщение 0 --- базовая инструкция роли; запрос пользователя --- первый ход
\texttt{user}. Карта состояний не входит в базовую инструкцию: инструкция запроса называет
стартовые состояния, инструкция перехода --- достижимые состояния и правила, локально для
текущего шага. Открытое состояние не видит цель другого состояния. Мотивация --- не
создавать «вторую цель» из карты переходов; это дизайн-гипотеза, не измеренная в работе.

\begin{figure}[t]
\centering
\begin{tikzpicture}[
  every node/.style={font=\small},
  box/.style={draw, rounded corners=2pt, align=center, inner sep=5pt, minimum height=8mm},
  arr/.style={-{Stealth[length=2mm]}},
  lbl/.style={font=\scriptsize},
]
\node[box] (journal) at (0,0) {журнал\\(только дописывание)};
\node[box] (proj) at (4.1,0) {$project()$};
\node[box] (ctx) at (7.6,0) {\texttt{Context}\\(вид)};
\node[box] (llm) at (11.2,0) {модель};
\node[box] (exec) at (7.6,-2.1) {исполнитель\\инструментов};
\draw[arr] (journal) -- (proj);
\draw[arr] (proj) -- (ctx);
\draw[arr] (ctx) -- (llm);
\draw[arr] (llm) |- node[lbl,right,pos=0.25]{действие} (exec);
\draw[arr] (exec) -| node[lbl,below,pos=0.25]{события (дописывание)} (journal);
\draw[arr,dashed] (llm.north) .. controls +(0,1.3) and +(0,1.3) ..
  node[lbl,above]{\texttt{recall} / \texttt{search}} (journal.north);
\end{tikzpicture}
\caption{Цикл работы: журнал --- единственный источник истины; контекст --- производный
вид, пересчитываемый каждый ход. Пунктир --- адресуемый доступ к скрытой истории.}
\label{fig:cycle}
\end{figure}

\subsection{Сессия и операции}

Агент сессии разбирает каждое сообщение пользователя и либо отвечает, либо открывает
единицу работы с самодостаточным заданием (табл.~\ref{tab:ops}). Каждая операция открывает
свой журнал: работа над документом не несёт чтений файлов из исследования, исправление
дефекта не тащит план.

\begin{table}[h]
\centering
\small
\caption{Операции сессии.}
\label{tab:ops}
\begin{tabular}{@{}ll@{}}
\toprule
Операция & Назначение \\
\midrule
\texttt{respond} & ответ из диалога и прошлой работы \\
\texttt{general} & исследование только на чтение, без изменений \\
\texttt{revision} & изменяющая операция, уже описанная в запросе \\
\texttt{task} & задача кодирования с переформулированным запросом \\
\texttt{plan} & план работы, только на чтение \\
\texttt{write\_document} & большой Markdown-документ по разделам \\
\bottomrule
\end{tabular}
\end{table}

\subsection{Состояния: контракт как данные}
\label{sec:states}

Состояние определяется четвёркой: цель (\texttt{what}), критерий остановки (\texttt{Done
when}), набор инструментов, разрешённые переходы. Цель и критерий --- данные журнала, попадающие
в рабочую инструкцию в точке решения. Переходы описываются протоколом: инструкция
перехода называет достижимые состояния и правила; решение принимает модель; исполнитель
записывает факт перехода в журнал.

Точнее говоря, это автомат, у которого \emph{структура} (узлы, контракты, допустимые
переходы) --- данные, а \emph{функция переходов} --- модель. Сейчас реализован один граф
внутри \texttt{task}; стартовать можно в Patch, Research или Construction, а переходы такие:
\[
\begin{aligned}
&\text{Research} \rightarrow \{\text{End},\ \text{Patch}\}, \\
&\text{Patch} \rightarrow \{\text{Verification},\ \text{Research}\}, \\
&\text{Construction} \rightarrow \{\text{End},\ \text{Research}\}, \\
&\text{Verification} \rightarrow \{\text{End},\ \text{Patch}\}.
\end{aligned}
\]
End терминален. Другие топологии архитектура допускает, но не реализует; делать из одного
графа общие выводы о топологиях преждевременно. Заметим: сам переход выбирает модель, поэтому
это не классический детерминированный конечный автомат, а автомат с топологией-данными и
моделью-решателем.

\begin{table}[h]
\centering
\small
\caption{Состояния единицы работы \texttt{task}: назначение и допустимые переходы.}
\label{tab:states}
\begin{tabular}{@{}lll@{}}
\toprule
Состояние & Назначение & Следующее \\
\midrule
Research & исследование рабочей области (чтение) & End, Patch \\
Patch & поиск и исправление дефекта & Verification, Research \\
Construction & создание нового артефакта по спецификации & End, Research \\
Verification & запуск проверок и тестов & End, Patch \\
End & завершение, итог родителю & --- \\
\bottomrule
\end{tabular}
\end{table}

Отличие от субагента --- не в наличии цели как таковой (у субагента цель есть --- текст
подзадачи), а в том, что цель, критерий и переходы живут в журнале, а факты переходов
журналируются и проверяемы постфактум. «Данные против кода» --- следствие, а не главный
тезис; отличие от \texttt{CompiledSubAgent} и LangGraph в том же: там контракт зашит в код
графа.

Практическое следствие, на которое мы рассчитываем: когда цель и критерий стоят в точке
решения, модель реже принимает запрос за более широкую задачу. Это гипотеза; как её проверить ---
в разделе~\ref{sec:future}. Оборотная сторона: пока критерий не проверяется машинно,
возможна преждевременная остановка --- мы наблюдали её (раздел~\ref{sec:discussion}).

\subsection{Динамические гипотезы: следующий шаг от состояний}
\label{sec:hypotheses}

В состояниях работа расписана заранее: фазы, их цели и переходы известны до начала работы, а
модель решает только, когда перейти дальше. Следующий шаг --- дать модели самой решать, из
каких шагов состоит работа. Для этого единица \texttt{task} умеет работать динамическими
гипотезами: модель выдвигает проверяемое утверждение, работает под ним, закрывает его и
переходит к следующему. Способ выбирается явно --- \texttt{states} или \texttt{hypotheses} ---
и запоминается, чтобы после паузы работа продолжилась так же.

Гипотеза --- это утверждение, которое может быть верным или неверным. Её открывает инструмент
\texttt{open\_hypothesis}. Модель вызывает его и пишет три поля: \texttt{label} (короткое имя
для панели), \texttt{claim} (что именно она собирается установить) и \texttt{done} (что нужно
увидеть, чтобы утверждение подтвердилось или опроверглось). Номер гипотезы --- \texttt{H1},
\texttt{H1.1} и так далее --- присваивает движок, модель его не пишет. Движок записывает
открытие в журнал. Пока гипотеза открыта, действия выполняются только под ней. Гипотеза может
быть простой: тогда под ней идут действия. Или составной: тогда она делится на подгипотезы.
Открытые гипотезы всегда образуют одну цепочку. Общий контейнер формулы $H$ создаёт движок;
модель его не пишет.

Закрывается гипотеза не вызовом инструмента, а текстовым ответом без вызовов --- JSON с
полями \texttt{verdict}, \texttt{observed} и \texttt{summary}. В нём модель пишет,
подтвердилось утверждение, опровергнуто или отброшено, что именно она увидела и каков итог.
Когда открытых гипотез не осталось, текст закрывает цель --- JSON с \texttt{outcome} и
окончательным ответом. Механизм здесь тот же, что у состояний: там модель вызывает
\texttt{move} и входит в состояние, а закрывает его текстом; здесь вместо \texttt{move} она
вызывает \texttt{open\_hypothesis}, а закрывает гипотезу таким же текстом. Ссылки
\texttt{follows\_from} показывают, на какие прошлые гипотезы опирается шаг; это только ссылка,
а не проверенный вывод.

Сообщения для модели собираются так же, как у состояний, но скрывают больше. Пока гипотеза
открыта, её работа видна. У закрытой остаются только строки открытия и закрытия, а все ходы
исчезают из сообщений и остаются в журнале, где их можно найти командами \texttt{recall} и
\texttt{search}. Поэтому чем больше гипотез закрыто, тем короче лента; журнал при этом хранит
всё.

Движок не проверяет ни утверждение, ни решение модели. Он записывает открытие и закрытие и
следит за структурой: поля не пустые, открыта одна цепочка, доступны нужные инструменты.
Условие проверки пишет модель, а не движок, поэтому её решение остаётся мнением о том, что она
увидела. У состояний ограничение то же (раздел~\ref{sec:limitations}), но у гипотезы оно
видно сразу: записаны и условие проверки, и наблюдение, поэтому закрытие можно проверить потом.

\section{Постановка эксперимента}

\subsection{Задачи и окружение}

13 задач, портированных из Terminal-Bench~\cite{terminalbench}. Каждая запускается в
Docker-контейнере: задача кладёт файлы, агент работает оболочкой и инструментами, затем
встроенная программа проверки выставляет оценку (0 или 1). Программы проверки создаются
заново из кэша задач в момент проверки, поэтому оценка не зависит от файлов, оставленных
агентом (ограничение доверия --- раздел~\ref{sec:limitations}).

Каждый прогон --- пара (задача, агент) при фиксированном сетевом режиме; $K=5$ повторов.
Сеть включается/выключается изолирующим прокси; при \texttt{off} сеть снимается только у
агента, а у программы проверки остаётся, потому что её скрипт сам ставит окружение. Фактор
проверяет, обходится ли агент без внешних ресурсов. Десять прогонов \texttt{crack-7z-hash}
при выключенной сети не запускались (задача помечена неподдерживаемой в этом режиме).
Оба агента используют одну и ту же модель \texttt{deepseek/deepseek-v4-flash-latest} с
низким уровнем рассуждений у одного поставщика, через один адрес API. Лимит --- число
ходов, без ограничения по времени; opencode --- версия 1.19.5. Набор задач и программы
проверки у обоих агентов одни и те же, поэтому по модели и проверкам сравнение честное;
различаются архитектура, промпты и инструменты (раздел~\ref{sec:limitations}).

\begin{table}[t]
\centering
\small
\caption{Протокол эксперимента.}
\label{tab:protocol}
\begin{tabular}{@{}ll@{}}
\toprule
Параметр & Значение \\
\midrule
Задачи & 13 портированных Terminal-Bench \\
Повторы & $K=5$ \\
Сетевые режимы & включена / выключена (изолирующий прокси) \\
Агенты & Skein; opencode 1.19.5 \\
Модель & \texttt{deepseek/deepseek-v4-flash-latest}, рассуждения: низкий; одна у обоих \\
Параллелизм & 10 (\texttt{fix-ocaml-gc} --- 3) \\
Лимит & 80 вызовов модели; \texttt{fix-ocaml-gc} --- 60; +5 в окне у End \\
Таймаут & нет \\
Прогоны & 260, запускаемых 250 \\
Дата серии & 2026-11-13 \\
\bottomrule
\end{tabular}
\end{table}

\subsection{Показатели и множества сравнения}

Фиксируются: оценка, число ходов, входные токены (новые и прочитанные из кэша), выходные
токены, токены рассуждений, стоимость. Доля решённых считается по всем запускаемым
прогонам. Токеновые сравнения --- по решённым прогонам (по 123 на агента). Множества решённых
\emph{различаются по составу} (провалы у агентов на разных задачах), поэтому агрегаты
содержат эффект отбора. Мы приводим обе версии: полные решённые множества
(табл.~\ref{tab:tokens}) и пересечение из 121 пары, где решены оба прогона
(табл.~\ref{tab:intersection}). Вывод на пересечении сохраняется.

opencode не возвращает стоимость. Тарифы получены точным решением линейной системы по
фактическому счёту поставщика для Skein (три коэффициента --- три агрегата счёта):
новый вход --- 17.80~руб./1M, чтение из кэша --- 0.356~руб./1M (в 50 раз дешевле нового
входа), выход --- 71.20~руб./1M. Запись в кэш в этих прогонах не тарифицировалась
(\texttt{cacheWrite = 0}); токены рассуждений входят в выход, поэтому отдельного тарифа на
них нет. Отношение 50$\times$ --- тарифное, а не измеренный эффект; измеренная доля кэша
во входе --- 85.7\% у Skein и 94.9\% у opencode. Для opencode рублёвая величина --- оценка
по той же модели, а не счёт. Для справки используем условный курс 80 руб./долл.: при нём
178.22 руб. --- это около 2.23 долл., а 125.65 руб. --- около 1.57 долл.

\section{Результаты}

\subsection{Решённые задачи: паритет на насыщенном наборе}

Из 250 запускаемых прогонов решено 246 (98.4\%); у каждого агента --- 123 из 125.
Запланировано 260 прогонов (13 задач $\times$ 2 агента $\times$ 2 сети $\times$ 5 повторов);
10 клеток \texttt{crack-7z-hash} при выключенной сети не запускались, поэтому исполнено 250,
а сравнения режимов --- на 12 общих задачах. По группам: Skein --- 63/65 (сеть вкл.) и 60/60 (выкл.); opencode --- 65/65 и 58/60.
Провалы Skein: \texttt{fix-git} (один повтор, преждевременная остановка) и
\texttt{fix-ocaml-gc} (один повтор, лимит ходов), оба при включённой сети. Провалы
opencode: два повтора \texttt{db-wal-recovery} при выключенной сети. Частоты: преждевременная
остановка --- 1 из 125 прогонов Skein (0.8\%), промах по лимиту --- 1 из 125 (0.8\%); у
opencode два провала из 125 (1.6\%), оба на одной задаче и режиме. По всем прогонам, включая
провалы, Skein стоит 182.72~руб., opencode --- \textasciitilde138.44~руб.; решённых у обоих
123 из 125. На уровне прогонов дискордантных пар 2 и 2 (точный тест МакНемара, p=1.0):
различий не обнаружено; Wilson 95\% для 123/125 --- [94.4\%; 99.6\%]. Это не доказательство
эквивалентности, а отсутствие обнаруженных различий.

Существенная оговорка: наш портированный поднабор насыщен (98.4\%) и не различает
агентов по качеству; сравнение содержательно только по структуре затрат. Внешние цифры по
полному Terminal-Bench мы не приводим как контраст: версии набора и лидерборд обновляются.
Для проверки заявленных выгод нужны длинные и трудные задачи (раздел~\ref{sec:future}).

\begin{figure}[t]
\centering
\begin{tikzpicture}
\begin{axis}[
  width=0.95\textwidth, height=6cm,
  ybar, bar width=14pt,
  ymin=0, ymax=70,
  ylabel={прогонов},
  symbolic x coords={skein_on,skein_off,opencode_on,opencode_off},
  xtick=data,
  xticklabels={Skein (вкл.),Skein (выкл.),opencode (вкл.),opencode (выкл.)},
  enlarge x limits=0.18,
  legend style={at={(0.5,-0.18)},anchor=north,legend columns=3,draw=none,font=\small},
  tick label style={font=\small},
  label style={font=\small},
]
\addplot[fill=skeinc,draw=none] table[x=group,y=solved,col sep=comma]{../figures/fig_solved.csv};
\addplot[fill=outc,draw=none] table[x=group,y=failed,col sep=comma]{../figures/fig_solved.csv};
\addplot[fill=cachec,draw=none] table[x=group,y=unsupported,col sep=comma]{../figures/fig_solved.csv};
\legend{решено,провал,не запускалось}
\end{axis}
\end{tikzpicture}
\caption{Исходы по группам. Всего 260 прогонов; 10 не запускались (\texttt{crack-7z-hash},
сеть выключена).}
\label{fig:solved}
\end{figure}

\subsection{Что показали журналы}

На 125 прогонах Skein инструмент \texttt{recall} не вызывался ни разу, \texttt{search} ---
два раза в двух прогонах; при этом \texttt{bash} вызывался 1985 раз, \texttt{read} --- 170,
\texttt{grep} --- 8. То есть адресуемый доступ к скрытой истории на этих коротких задачах
почти не использовался, и выгоды \texttt{recall}/\texttt{search} здесь не проверяются.
Состояния использовались умеренно: 301 событие \texttt{move} на 125 прогонов (медиана 2);
движок отклонил 22 действия из 2518 (0.9\%).

\subsection{Токены и стоимость: меньше объём, дороже счёт}

opencode тратит в 1.49 раза больше токенов суммарно, но почти весь избыток --- чтение из
кэша (94.9\% входа против 85.7\% у Skein). Skein компактнее по суммарному входу, однако
новых входных токенов у него в 1.87 раза больше, а выходных --- в 1.33 раза больше:
единицы работы не делят общего растущего начала и могут читать рабочую область заново.
Ходов у Skein также больше (в 1.36 раза). При 50-кратной разнице тарифов на кэш и новый
вход это делает opencode дешевле: 125.65 против 178.22~руб., то есть opencode дешевле на 29.5\%, Skein дороже в
1.42 раза (табл.~\ref{tab:tokens}, рис.~\ref{fig:tokens}).

Доля кэша в первом повторе близка к последующим (Skein 86.7\% против 85.3\%; opencode
94.4\% против 95.2\%), поэтому межпрогонное переиспользование префикса не объясняет разрыв.
На уровне задач opencode дешевле в 12 из 13 (знаковый тест, p$\approx$0.0034); дороже только
на \texttt{fix-ocaml-gc}.

Разложение: по шагам ходов у Skein в 1.36 раза больше, по API-вызовам --- в 1.27; выход на
ход почти равен (507 против 520), а новый вход на API-вызов больше в 1.47 раза (на шаг --- в
1.37). End даёт 9.8\% ходов; без End отношение шагов 1.22. Значит, рост нового входа --- это
большее число вызовов при почти том же выходе на вызов, а не более длинные ответы.

Счётчики инструментов за 125 прогонов: Skein --- 1985 \texttt{bash}, 170 \texttt{read},
113 \texttt{write}, 70 \texttt{edit}, 48 \texttt{list}, 8 \texttt{grep}; opencode ---
1982 \texttt{bash}, 75 \texttt{read}, 83 \texttt{write}, 70 \texttt{edit}
(\texttt{grep}/\texttt{list} у него нет — поиск идёт через \texttt{bash}). Явных
\texttt{read} у Skein больше (170 против 75), а \texttt{bash} примерно столько же; это не
доказывает повторное чтение между состояниями (раздел~\ref{sec:discussion}).

По API-вызовам: Skein --- 2447 (один ход = один вызов), opencode --- 1923 при 1800 шагах
(1.07 вызова на шаг); на решённых прогонах отношение 1.27, а не 1.36 по шагам. Новый вход на
API-вызов --- 1840 против 1252 токенов (×1.47; на шаг --- ×1.37). Число \texttt{bash}-вызовов
почти равно (1985 против 1982), поэтому разница в ходах --- это гранулярность и протокол, а
не больший объём терминальной работы. Ход Skein --- один вызов модели; в нём может быть
несколько действий, поэтому счётчики действий превышают число ходов, а \texttt{move} ---
действие внутри хода, а не отдельный ход.

\begin{table}[t]
\centering
\small
\caption{Токены и стоимость, только решённые прогоны (по 123 на агента; состав множеств
различается, см.~\S4.2). Отношения посчитаны по неокруглённым значениям. Ходы: у Skein ---
вызовы модели, у opencode --- шаги агента; API-вызовы модели у opencode --- отдельная
строка. Стоимость opencode --- оценка.}
\label{tab:tokens}
\begin{tabular}{@{}lrrr@{}}
\toprule
 & Skein & opencode & opencode/Skein \\
\midrule
Вход, всего (млн) & 31.38 & 47.63 & 1.52 \\
\quad из кэша (млн) & 26.88 & 45.22 & 1.68 \\
\quad новый (млн) & 4.50 & 2.41 & 0.53 \\
Выход (млн) & 1.24 & 0.94 & 0.75 \\
Рассуждения (млн) & 0.81 & 0.62 & 0.77 \\
Всего, вход+выход (млн) & 32.62 & 48.56 & 1.49 \\
Доля кэша во входе & 85.7\% & 94.9\% & --- \\
Ходы/шаги & 2447 & 1800 & 0.74 \\
API-вызовы модели & 2447 & 1923 & 0.79 \\
Стоимость (руб.) & 178.22 & \textasciitilde125.65 & 0.71 \\
\bottomrule
\end{tabular}
\end{table}

\begin{table}[t]
\centering
\small
\caption{Проверка на пересечении: 121 пара прогонов, где решены оба агента. Стоимость
opencode --- оценка. Отношения --- opencode/Skein.}
\label{tab:intersection}
\begin{tabular}{@{}lrrr@{}}
\toprule
 & Skein & opencode & opencode/Skein \\
\midrule
Вход, всего (млн) & 31.08 & 45.11 & 1.45 \\
\quad из кэша (млн) & 26.64 & 42.79 & 1.61 \\
\quad новый (млн) & 4.44 & 2.32 & 0.52 \\
Выход (млн) & 1.22 & 0.92 & 0.76 \\
Рассуждения (млн) & 0.80 & 0.62 & 0.78 \\
Ходы/шаги & 2407 & 1751 & 0.73 \\
Стоимость (руб.) & 175.60 & 122.26 & 0.70 \\
\bottomrule
\end{tabular}
\end{table}

На пересечении Skein дороже в 1.44 раза (opencode дешевле на 30.4\%), то есть вывод не
меняется. Парный бутстрэп по 121 паре (5000 пересэмплирований) даёт 95\%-й интервал для
отношения стоимости opencode/Skein 0.70 [0.59--0.81]; кластерный бутстрэп по 13 задачам ---
[0.50--0.85]. Для отношения нового входа интервал --- 0.52 [0.44--0.61]. Все три интервала не
включают 1, то есть направление вывода устойчиво к отбору и кластеризации. Спаривание по
номеру повтора искусственно (прогоны независимы), а покрытие кластерного интервала по 13
задачам нестабильно; поэтому интервалы --- описательные, а не строгий парный дизайн.

\begin{figure}[t]
\centering
\begin{tikzpicture}
\begin{axis}[
  width=0.8\textwidth, height=6cm,
  ybar stacked, bar width=26pt,
  ymin=0, ymax=55,
  ylabel={млн токенов},
  symbolic x coords={Skein,opencode},
  xtick=data,
  enlarge x limits=0.5,
  legend style={at={(0.5,-0.16)},anchor=north,legend columns=3,draw=none,font=\small},
  tick label style={font=\small},
  label style={font=\small},
]
\addplot[fill=freshc,draw=none] table[x=agent,y=fresh,col sep=comma]{../figures/fig_tokens.csv};
\addplot[fill=cachec,draw=none] table[x=agent,y=cache,col sep=comma]{../figures/fig_tokens.csv};
\addplot[fill=outc,draw=none] table[x=agent,y=output,col sep=comma]{../figures/fig_tokens.csv};
\legend{новый вход,кэш,выход}
\end{axis}
\end{tikzpicture}
\caption{Состав токенов, решённые прогоны (агрегаты по 123 прогонам на агента). Избыток
объёма у opencode почти целиком --- кэш; у Skein выше доля нового входа.}
\label{fig:tokens}
\end{figure}

Медианы и межквартильный размах на прогон приведены в табл.~\ref{tab:medians}. Медианный
прогон opencode дешевле; Skein не стремится к минимуму стоимости, поэтому медиана денег
мало говорит о качестве управления.

\begin{table}[t]
\centering
\small
\caption{Медианы на прогон и межквартильный размах (в скобках) на 12 общих задачах (без
\texttt{crack-7z-hash}). Ходы: у Skein --- вызовы модели, у opencode --- шаги. Стоимость
opencode --- оценка.}
\label{tab:medians}
\begin{tabular}{@{}lrrrr@{}}
\toprule
Агент / сеть & Ходы/шаги & Новый вход, тыс. & Выход, тыс. & Стоимость, руб. \\
\midrule
Skein, вкл. & 14 [11--30] & 22.2 [17.7--56.8] & 4.5 [2.8--14.5] & 0.70 [0.54--2.25] \\
Skein, выкл. & 15 [10.8--24.3] & 25.0 [17.4--46.7] & 7.3 [3.0--12.3] & 1.00 [0.58--1.99] \\
opencode, вкл. & 9 [7--20] & 10.3 [4.1--26.5] & 2.9 [1.6--10.1] & 0.40 [0.20--1.42] \\
opencode, выкл. & 9 [7--19] & 10.4 [4.7--26.4] & 2.7 [1.5--11.4] & 0.44 [0.23--1.37] \\
\bottomrule
\end{tabular}
\end{table}

На 12 общих задачах аномалия сохраняется: медианный выход Skein в \texttt{off} выше, чем в
\texttt{on} (7.3k против 4.5k), у opencode ниже (2.7k против 2.9k); это не эффект состава
задач. Разница 1.42$\times$ --- по суммам; медианы дают 0.70 против 0.40 (отношение 1.75),
но IQR сильно перекрываются, а Skein дешевле в 23 из 121 пары (19\%). Направление устойчиво,
но величина зависит от агрегата.

\subsection{Декомпозиция по состояниям Skein}

По всем 125 прогонам Skein дороже всего Construction (35.5\% стоимости) и Patch (26.1\%),
затем Research (23.8\%). Verification дешёва: 8.6\% стоимости при 11.6\% ходов. Patch и
Verification вместе --- 34.7\% стоимости и 36.7\% ходов (925 из 2518). End в среднем
занимает около двух ходов на прогон (248/125) и 6.0\% стоимости --- накладные расходы
завершения. Construction дорог в основном за счёт \texttt{cobol-modernization} (32.7~руб.)
и \texttt{regex-log} (14.0~руб.) --- задач создания артефакта с нуля; это состав задач, а не
неэффективность самого состояния. Разбивка приблизительна: метку состояния выбирает модель, ход перехода
отнесён к новому состоянию.

\begin{table}[t]
\centering
\small
\caption{Стоимость по состояниям Skein (все 125 прогонов, включая провалы).}
\label{tab:state}
\begin{tabular}{@{}lrrr@{}}
\toprule
Состояние & Ходы & Доля ходов & Стоимость \\
\midrule
Construction & 760 & 30.2\% & 64.86~руб. (35.5\%) \\
Patch & 634 & 25.2\% & 47.67~руб. (26.1\%) \\
Research & 585 & 23.2\% & 43.49~руб. (23.8\%) \\
Verification & 291 & 11.6\% & 15.71~руб. (8.6\%) \\
End & 248 & 9.8\% & 10.98~руб. (6.0\%) \\
\bottomrule
\end{tabular}
\end{table}


\subsection{Гипотезы против состояний}
\label{sec:hyp-exp}

Гипотезы измеряются на семи задачах, где проверяемое утверждение естественно: поиск причины и
восстановление (\texttt{fix-git}, \texttt{git-leak-recovery},
\texttt{custom-memory-heap-crash}, \texttt{password-recovery}, \texttt{db-wal-recovery},
\texttt{fix-code-vulnerability}, \texttt{fix-ocaml-gc}). Условия те же: $K=5$, сеть включена и
выключена, одна модель. Оба способа работы гоняются в одной серии, поэтому решённые задачи,
ходы и новый вход сравнимы; рублёвые стоимости сравнимы только внутри этой серии, потому что
тариф дня отличается от серии с состояниями. Лимит --- штатный лимит задачи; на
\texttt{fix-ocaml-gc} гипотезы дополнительно прогнаны с лимитом 80 ходов. Проверка на задаче
с нуля --- \texttt{modernize-scientific-stack}.

\begin{table}[t]
\centering
\small
\caption{Состояния и гипотезы на семи задачах ($K=5$, сеть включена/выключена). Медианы на
прогон; стоимость --- по тарифу дня серии, сравнивать только внутри таблицы.}
\label{tab:hyp}
\begin{tabular}{@{}lrrrr@{}}
\toprule
Способ / сеть & Решено & Ходы & Новый вход, тыс. & Стоимость, руб. \\
\midrule
Состояния, вкл. & 35/35 & 17 & 34.5 & 2.42 \\
Состояния, выкл. & 34/35 & 15 & 28.7 & 1.61 \\
Гипотезы, вкл. & 34/35 & 20 & 19.6 & 1.88 \\
Гипотезы, выкл. & 34/35 & 22 & 20.1 & 1.98 \\
\bottomrule
\end{tabular}
\end{table}

Гипотезы решают столько же задач, сколько состояния: 34--35 из 35 в каждой группе.
Единственный провал --- не хватило ходов на длинной задаче. Основное различие --- в сообщениях,
которые видит модель: гипотезы тратят почти вдвое меньше новых входных токенов (19.6--20.1 тыс.
против 28.7--34.5 тыс. на запуск), потому что работа закрытой гипотезы исчезает из ленты. За это
приходится платить ходами: 20--22 против 15--17, открытие и закрытие гипотезы стоят отдельных
шагов. По деньгам картина смешанная: на коротких задачах гипотезы дешевле (\texttt{fix-git}
0.79/0.98 против 1.43/1.30~руб., \texttt{custom-memory} 3.2/3.5 против 3.9/5.0,
\texttt{fix-code} 1.6/0.83 против 2.9/1.6), а на \texttt{db-wal-recovery} без сети и на
длинной задаче дороже по ходам.

\texttt{fix-ocaml-gc} показывает границу способа. При лимите 60 гипотезы решают 8 из 10 (оба
провала --- \texttt{max\_turns}), состояния --- 9 из 10. При лимите 80 гипотезы решают 10 из
10 (медианы 44 и 68 ходов, 1.98 и 3.42~руб.). Значит, провал при 60 --- следствие бюджета, а
не способа: длинной задаче нужен запас ходов.

Контрольная задача \texttt{modernize-scientific-stack}: гипотезы решают 9 из 10 (18 ходов,
около 0.58~руб.). Единственный провал показателен. Скрипт посчитал среднее верно
($-15.5^\circ$C), но напечатал перед ним лишнюю строку с диапазоном; проверка ищет первое число
после «Station 101» и берёт $-4.3$. Расчёт верный, но ответ не приняли. Это не свойство
гипотез, а общее ограничение (раздел~\ref{sec:limitations}): модель не решает, принят ли
результат, а проверка смотрит и на формат. С гипотезами это заметнее: модель записала, что цель
достигнута, а принимал результат не она.

\subsection{Что показывают журналы гипотез}
\label{sec:hyp-journals}

Гипотезы не остаются декларацией: они действительно создаются и закрываются. В 70 прогонах
гипотезного способа хотя бы одна гипотеза есть в каждом прогоне; всего их 99, чаще всего одна
на прогон (47 прогонов), две --- в 19, три-четыре --- в четырёх. Закрыты 95: 93 подтверждены,
2 опровергнуты, 0 отброшены. Ещё 4 остались открытыми, и все они --- в четырёх прогонах
\texttt{fix-ocaml-gc}, которые упёрлись в лимит 60 ходов; два из этих прогонов всё равно прошли
проверку. Вложенные гипотезы (\texttt{H1.1}) редки: всего 4; ссылки \texttt{follows\_from}
заполнены у 21 гипотезы из 99.

Гипотезы отражают ход решения. В каждом из 66 прогонов, где цель закрыта, она закрыта после
последней закрытой гипотезы: пока есть открытая, работа не считается законченной. Чем больше
гипотез, тем больше ходов: в прогонах с одной гипотезой медиана 20 ходов, с двумя и более ---
28. На длинной \texttt{fix-ocaml-gc} рост бюджета с 60 до 80 ходов увеличивает и число гипотез
(16 $\rightarrow$ 21), и число закрытых целей (6 $\rightarrow$ 9 из 10): дополнительные ходы
уходят на новые проверяемые шаги, а не на повторение сделанного.

При этом закрытая цель не означает, что результат принят. В контрольной задаче
\texttt{modernize} все 10 прогонов закрыли цель как \texttt{achieved}, а проверка приняла 9: в
одном случае расчёт был верным, но формат ответа --- нет. Журнал показывает, что модель считает
работу законченной; принимает результат по-прежнему внешняя проверка.

\section{Обсуждение}
\label{sec:discussion}

\subsection{Экономика кэша и выхода}

Тариф на чтение кэша в 50 раз ниже тарифа нового входа (17.80 / 0.356), а измеренная доля
кэша во входе --- 85.7\% у Skein и 94.9\% у opencode. Поэтому агент с большим суммарным
объёмом (opencode) оказывается дешевле агента с малым объёмом, но большей долей нового
входа (Skein).

Разрыв 52.6~руб. раскладывается не только на кэш (табл.~\ref{tab:costsplit}): новый вход
даёт +37.2~руб., выход --- +21.6~руб., а кэш экономит 6.5~руб. Выход --- крупнейшая статья
стоимости у обоих агентов (49.5\% у Skein, 53.1\% у opencode), и около 66\% выхода ---
рассуждения. Без кэша Skein дешевле на 29\% (646.9 против 914.5~руб.): кэш переворачивает
знак (при текущем тарифе и текущем cache hit rate). Если бы при своём входе Skein имел долю
кэша opencode (94.9\%), стоимости почти сравнялись бы (127.4 против 125.7~руб.). Отсюда
цель оптимизации формулируется точнее: стабилизация префикса внутри состояния, а не
сокращение числа токенов само по себе.

Чувствительность к тарифу (табл.~\ref{tab:sensitivity}): при отношении fresh/cache 5$\times$
Skein немного дешевле (264.1 против 270.6~руб.), при 10$\times$ opencode дешевле на 12\%,
при 50$\times$ --- на 29\%, при 100$\times$ --- на 32\%. Break-even --- около 5.5$\times$.
Чтобы Skein сравнялся при текущем тарифе, ему нужна доля кэша 95.2\% (сейчас 85.7\%).

\begin{table}[t]
\centering
\small
\caption{Стоимость при разных отношениях тарифов fresh/cache (решённые прогоны; объёмы
фиксированы по измеренным, пересчитываются только тарифы). При 5$\times$ Skein немного
дешевле; break-even --- около 5.5$\times$.}
\label{tab:sensitivity}
\begin{tabular}{@{}rrrr@{}}
\toprule
fresh/cache & Skein, руб. & opencode, руб. & opencode/Skein \\
\midrule
5   & 264.1 & 270.6 & 1.03 \\
10  & 216.2 & 190.1 & 0.88 \\
20  & 192.3 & 149.8 & 0.78 \\
50  & 178.0 & 125.7 & 0.71 \\
100 & 173.2 & 117.6 & 0.68 \\
\bottomrule
\end{tabular}
\end{table}

Практический вывод: при действующем кэшировании суммарные токены --- плохой прокси
стоимости, а стоимость --- плохой прокси качества управления контекстом. Сравнивать
архитектуры по деньгам на коротких задачах малоинформативно в обе стороны.

\begin{table}[t]
\centering
\small
\caption{Декомпозиция разрыва стоимости (Skein $-$ opencode, решённые прогоны). Сумма
отличается от 52.6~руб. из-за округления.}
\label{tab:costsplit}
\begin{tabular}{@{}lrr@{}}
\toprule
Компонент & Вклад, руб. & Доля разрыва \\
\midrule
Новый вход ($+2.09$M $\times$ 17.80) & $+37.2$ & $+71\%$ \\
Выход ($+0.30$M $\times$ 71.20) & $+21.6$ & $+41\%$ \\
Кэш ($-18.34$M $\times$ 0.356) & $-6.5$ & $-12\%$ \\
\midrule
Итого & $+52.3$ & --- \\
\bottomrule
\end{tabular}
\end{table}

\subsection{Цена текущей реализации: префикс и повторное чтение}

Изоляция контекста не бесплатна: краткий итог предыдущей единицы не несёт прочитанного
содержимого, поэтому следующая единица может читать рабочую область заново; в ленте факты
переиспользуются из кэша. Измерение возможно только для явных \texttt{read}: в журналах
Skein их 170, из них 9 (5.3\%) --- файл, уже прочитанный в другом состоянии; чтения через
\texttt{bash} (1985 вызовов) не выделены. Первый рабочий ход состояния имеет cache hit
66.2\% против 87.1\% на последующих ходах; это прямое свидетельство нестабильности префикса
на входе в состояние, хотя первые ходы дают лишь около 2\% входа. Эффект входа в
состояние объясняет лишь около 0.4 п.п. из 9.2 п.п. разрыва доли кэша; остальное ---
поведение префикса внутри состояния. Мы интерпретируем рост
нового входа как цену текущей реализации изоляции --- нестабильность префикса и большее
число ходов, --- а не как доказанное свойство архитектуры. Иллюстративный случай (не часть эксперимента:
другая постановка, один репозиторий, шесть реплик, $K=1$): Skein --- 63 шага против 24 у
opencode, 113.7k нового входа против 35.8k, 3.71~руб. против \textasciitilde1.46~руб.,
доля кэша 67\% против 95\% на короткой сессии. Возможен ли выигрыш на длинной работе ---
открытый вопрос. Предсказание: у монолитной ленты сжатие или усечение может инвалидировать
префикс, и тогда преимущество кэша opencode исчезает именно на длинных задачах; это
гипотеза, а не измерение. Две самые длинные по сумме ходов задачи --- \texttt{fix-ocaml-gc} (392) и
\texttt{custom-memory-heap-crash} (390). На первой Skein дешевле в расчёте на решённый
прогон (2.58 против 2.81~руб.), на второй --- дороже (2.87 против 1.38~руб.); это скорее
эффект отбора, чем свойство архитектуры.

\subsection{Состояния и субагенты: что именно отличается}

Функционально \texttt{general}, \texttt{revision}, \texttt{plan}, \texttt{write\_document}
--- субагенты. Содержательное отличие \texttt{task} --- контракт как данные: цель, критерий
и топология переходов живут в журнале и интерпретируются протоколом, факты переходов
записываются. Это делает топологии предметом исследования (их можно менять, сравнивать и
воспроизводить без изменения кода), а не свойством реализации. Пока реализован один
граф состояний; вопрос о других топологиях (миграции, расследования инцидентов, документ с
проверкой фактов) мы оставляем открытым --- это отдельное направление, не текущая фаза проекта.

\subsection{Состояния и гипотезы: кто разбивает работу на шаги}

Оба способа используют один журнал и одни и те же сообщения. Разница в том, кто разбивает
работу на шаги. Состояния задают это заранее: фазы, их цели и переходы известны до начала
работы, модель решает только, когда перейти дальше. В гипотезах модель решает это сама по ходу
работы: она выдвигает проверяемые утверждения, и в сообщениях видна только работа открытой
гипотезы. Поэтому состояния лучше подходят для задач, где есть готовая спецификация или
процедура, а гипотезы --- для задач, где нужно искать причину и что-то проверять. Проверка на
задаче с нуля подтверждает это лишь отчасти: гипотезы решают её, но тратят больше ходов и
чувствительны к формату результата.

Общее у обоих способов --- условие проверки в момент решения. У состояния это
\texttt{Done when}, у гипотезы --- \texttt{done} и записанное \texttt{observed}. В обоих
случаях движок не проверяет условие, а записывает мнение модели. Поэтому решение модели не
равно истине: оно позволяет проверить шаг потом, но не доказывает его. Разница в том, что у
гипотезы условие и наблюдение есть на каждом шаге, а у состояния --- только на границе.

\subsection{Провалы}

\texttt{fix-git}, Skein, сеть вкл., r4: агент нашёл висячий коммит, но не слил его с
веткой; состояние Research закрылось текстовым диагнозом, после чего модель перешла в End.
Переход Research$\rightarrow$End топологией разрешён, но его предусловие --- «требуемый
результат существует» --- не выполнено. Движок не проверяет критерий, поэтому переход
записан как допустимый. Это преждевременная остановка, а не недопустимый переход. В
остальных повторах слияние выполнено.

\texttt{fix-ocaml-gc}, Skein, сеть вкл., r5: лимит 60 ходов исчерпан; при приближении к End
движок дал окно ещё в 5 ходов, поэтому прогон остановился на 65. Программа проверки не нашла
ожидаемого вывода тестов. Промах по лимиту, не сбой стенда.

\texttt{db-wal-recovery}, opencode, сеть выкл., r2 и r4: 40--57 шагов, вход 2.3M--4.2M
токенов (в основном кэш). Причина не разобрана; материал для анализа, не вывод о методе.

\section{Ограничения}
\label{sec:limitations}

\begin{itemize}[leftmargin=1.4em,itemsep=0.2em]
  \item \textbf{Насыщенный набор.} 98.4\% решённых: поднабор не различает качество;
  выводы ограничены структурой затрат на коротких задачах.
  \item \textbf{Одна модель, один поставщик API.} Низкий уровень рассуждений; переносимость
  не проверялась. Точный snapshot модели не зафиксирован: в имени стоит \texttt{latest}.
  \item \textbf{Разные инструменты и промпты.} Наборы инструментов у агентов различаются;
  абляция \texttt{recall}/\texttt{search} не проводилась, поэтому эффект архитектуры не
  отделён от эффекта инструментального интерфейса.
  \item \textbf{Стоимость opencode --- оценка} по тарифам, выведенным из счёта Skein.
  Токеновые выводы надёжнее денежных.
  \item \textbf{Разный состав решённых множеств.} Агрегаты по 123 решённым прогонам
  содержат эффект отбора; пересечение --- 121 пара, результаты на нём приведены в
  табл.~\ref{tab:intersection}.
  \item \textbf{Малая выборка.} $K=5$, 13 задач; медианы и межквартильный размах приведены,
  бутстрэп-интервалы для отношений посчитаны, но полноценных кластерных тестов нет.
  \item \textbf{Загрязнение.} Задачи Terminal-Bench публичны и могли попасть в обучение
  модели; мы это не контролировали.
  \item \textbf{Критерий остановки не проверяется машинно.} Состояние закрывается
  текстовым ответом модели; обнаружение циклов и ранней остановки отложено.
  \item \textbf{Детерминированность проекции $\neq$ воспроизводимость поведения.}
  Проекция проверена на 400 синтетических и 125 реальных журналах (идентичный контекст);
  поведение агента стохастично.
  \item \textbf{Ключевые выгоды не измерены.} То, что агент удерживает цель и им можно
  управлять, --- пока предположение, а не результат.
  \item \textbf{Повторные чтения не измерены напрямую.} В журналах Skein \texttt{read} ---
  170 вызовов, 9 из них (5.3\%) повторяют файл в другом состоянии; основная работа идёт
  через \texttt{bash}, поэтому цену повторного чтения мы не выделяем.
  \item \textbf{Одна топология состояний}; учёт по состояниям приблизителен (метку
  выбирает модель).
  \item \textbf{Доверие к проверке.} Программы проверки перекладываются после агента, но
  системные бинарники агент теоретически мог изменить; изоляция проверки не усилена.
  \item \textbf{Prompt injection.} \texttt{recall}/\texttt{search} возвращают недоверенный
  текст файлов в ленту; для агента на чужих репозиториях это поверхность атаки, мы её не
  защищали. Для чужих документов (например, договоров) это ещё и риск принять текст за
  инструкцию.
  \item \textbf{Параметры сэмплирования не зафиксированы.} Температура и связанные
  параметры в протоколе не указаны.
  \item \textbf{Отбор задач.} 13 задач --- доступные портированные задачи Terminal-Bench,
  а не случайная выборка; внешняя валидность ограничена.
  \item \textbf{Стоимость проекции.} \texttt{project()} пересчитывает контекст по всему
  журналу за ход ($O(n)$ на ход); для длинных задач нужны снапшоты или мемоизация.
  \item \textbf{Взаимодействие с режимом сети не объяснено.} Медианный выход Skein в
  \texttt{off} выше, чем в \texttt{on} (7.3k против 4.5k), у opencode ниже; мы это не
  интерпретируем.
  \item \textbf{Окружение.} Только \texttt{/bin/bash} (Windows --- WSL2/Docker);
  13 портированных задач, не полный набор.
  \item \textbf{Решение по гипотезе не проверяется.} Движок проверяет только структуру закрытия
  (поля не пустые, открыта одна цепочка); утверждение и решение --- мнение модели.
  \item \textbf{Серия с гипотезами меньше.} Семь задач, $K=5$, одна модель; проверка на задаче
  с нуля --- одна задача. Рублёвые стоимости гипотез сравнимы только внутри этой серии.
  \item \textbf{Способ выбирается явно.} Система не решает сама, где нужны гипотезы, а где
  состояния; выбирает пользователь.
  \item \textbf{\texttt{follows\_from} не проверяется.} Ссылки сохраняются, но движок не
  проверяет, что вывод действительно следует из указанных гипотез.
  \item \textbf{Длинной задаче нужно больше ходов.} На \texttt{fix-ocaml-gc} гипотезам нужен
  бюджет 80 ходов; при 60 они упираются в лимит, хотя задача решаема.
  \item \textbf{Проверка смотрит и на формат.} На \texttt{modernize} верный расчёт не прошёл
  из-за лишней строки в выводе: проверка чувствительна к формату, и решение модели этого не
  меняет.
\end{itemize}

\section{Доступность}

Исходный код, документация и материалы кампании:
\url{https://github.com/merfill/skein}. Описание проекта и установка:
\url{https://skein-swe.dev}. Пакет npm \texttt{skein-swe}:
\url{https://www.npmjs.com/package/skein-swe}; образ Docker:
\url{https://hub.docker.com/r/mefrill/skein-swe} (аккаунт Docker Hub \texttt{mefrill}, не
\texttt{merfill}). Документация --- в \texttt{docs/en} и \texttt{docs/ru}; материалы
кампании --- в \texttt{bench/results/campaign-k5/}; журналы клеток --- в \texttt{result.json}
(поле \texttt{events}), replay 125/125; сырые прогоны публикуются с выпуском.

\section{Заключение и дальнейшая работа}
\label{sec:future}

Измерено: на насыщенном наборе коротких задач Skein и opencode решают одинаково; Skein
тратит меньше токенов суммарно, но больше нового входа и выхода, и при 50-кратной разнице
тарифов на кэш оказывается в 1.42 раза дороже. Мы интерпретируем это как цену текущей
реализации изоляции: нестабильность префикса и большее число ходов; при доле кэша, как у
opencode, стоимости почти равны (127.4 против 125.7 руб.). Дрейф внимания не измерялся; явные
повторные \texttt{read} редки (9 из 170), чтения через \texttt{bash} не выделены. Мы пока не
измеряли, удерживает ли агент цель в момент решения и насколько им можно управлять.

Следующая фаза --- развитие журнала как основы контекста, а не расширение функциональности
ради неё самой:
\begin{enumerate}[leftmargin=1.4em,itemsep=0.2em]
  \item немонотонный контекст: первый шаг уже сделан --- динамические гипотезы, работа
  закрытой гипотезы уходит с ленты и остаётся в журнале; дальше --- автоматически выбирать
  способ и проверять условия машинно: модель формулирует \texttt{done} как формальную
  спецификацию (предикат, регулярное выражение, шаблон AST), а движок проверяет её до закрытия;
  так же проверять ссылки \texttt{follows\_from};
  \item рассуждения (проект Ankyra --- исследовательский проект, где журнал используется как
  носитель цепочек рассуждений): журнал должен хранить не только действия и наблюдения, но и
  цепочки рассуждений, а проекция --- выбирать, что из них показывать;
  \item откаты до контрольных точек: журнал позволяет вернуться к сохранённому состоянию
  работы и продолжить с другой ветки; это операция над журналом, а не над лентой;
  \item агент наблюдения: сейчас он нерабочий; его восстановление --- снова работа с
  журналом, потому что наблюдатель читает журнал и предлагает подсказки или остановку.
  \item полная кампания гипотез на 13 задачах и вторая модель: текущая серия --- семь задач
  и одна модель; конструктивный домен измерен одной задачей.
\end{enumerate}

Архитектура с неизменяемым журналом и проверяемыми шагами особенно уместна там, где нужен
аудиторский след: в юридических и формальных текстах, где важно не только получить ответ, но
и показать, откуда он взялся.

\appendix

\section{Воспроизведение}

\begin{sloppypar}
Материалы: \texttt{bench/results/campaign-k5/} --- протокол (\texttt{manifest.json}),
сводки (\texttt{summary.md}, \texttt{summary.json}, \texttt{summary.csv}), сравнение
токенов (\texttt{token-comparison.md}), стоимость по состояниям (\texttt{state-cost.md}).
Сырые прогоны и журналы --- \texttt{bench/runs/matrix/campaign-k5/}; журналы клеток лежат в
\texttt{*/*cells/*/result.json} (поле \texttt{events}), replay 125/125.
\end{sloppypar}

Серия запускается через \texttt{tests/sandbox/campaign.ts}: \texttt{--repeat 5},
\texttt{--concurrency 10}, \texttt{--task-concurrency fix-ocaml-gc=3}, оба сетевых режима,
без ограничения по времени; пробный прогон без контейнеров --- \texttt{--dry}. Требуются
Docker, кэш задач, исполняемый файл opencode; \texttt{SKEIN\_LIVE=true}.

Серия с гипотезами: те же \texttt{campaign.ts}/\texttt{matrix.ts} с
\texttt{--agents skein,skein-hypotheses}; режим задаётся \texttt{--mode hypotheses} (стенд)
или \texttt{SKEIN\_TASK\_MODE=hypotheses} (агент). Артефакты ---
\texttt{bench/runs/matrix/inquiry-v26-modes},
\texttt{inquiry-v27-ocaml-80}, \texttt{inquiry-v28-modernize}.

\section{Результаты по задачам}

\begin{table}[h]
\centering
\small
\caption{Решённые прогоны и стоимость по задачам (только решённые; стоимость opencode ---
оценка). Знаменатели различаются, поэтому добавлены нормированные колонки руб./прогон.
У \texttt{crack-7z-hash} знаменатель 5 (сеть выкл. не запускалась). Суммы могут отличаться
от табл.~\ref{tab:tokens} на копейки из-за округления.}
\label{tab:pertask}
\begin{tabular}{@{}lrrrrrr@{}}
\toprule
Задача & Skein & opencode & S, руб. & O, руб. & S, руб./пр. & O, руб./пр. \\
\midrule
cobol-modernization & 10/10 & 10/10 & 36.73 & 34.09 & 3.67 & 3.41 \\
crack-7z-hash & 5/5 & 5/5 & 5.32 & 4.98 & 1.06 & 1.00 \\
custom-memory-heap-crash & 10/10 & 10/10 & 28.68 & 13.78 & 2.87 & 1.38 \\
db-wal-recovery & 10/10 & 8/10 & 8.64 & 5.95 & 0.86 & 0.74 \\
fix-code-vulnerability & 10/10 & 10/10 & 10.78 & 4.26 & 1.08 & 0.43 \\
fix-git & 9/10 & 10/10 & 5.44 & 1.95 & 0.60 & 0.20 \\
fix-ocaml-gc & 9/10 & 10/10 & 23.21 & 28.09 & 2.58 & 2.81 \\
git-leak-recovery & 10/10 & 10/10 & 8.27 & 1.82 & 0.83 & 0.18 \\
log-summary-date-ranges & 10/10 & 10/10 & 5.03 & 3.24 & 0.50 & 0.32 \\
modernize-scientific-stack & 10/10 & 10/10 & 4.98 & 2.33 & 0.50 & 0.23 \\
openssl-selfsigned-cert & 10/10 & 10/10 & 5.50 & 2.35 & 0.55 & 0.24 \\
password-recovery & 10/10 & 10/10 & 19.66 & 10.46 & 1.97 & 1.05 \\
regex-log & 10/10 & 10/10 & 15.99 & 12.35 & 1.60 & 1.23 \\
\bottomrule
\end{tabular}
\end{table}

\clearpage
\bibliographystyle{plainnat}
\bibliography{../refs}

\end{document}
